\section{Evaluation}\label{sec:eval}

This section reports measurements for xlog's deterministic, probabilistic, and neural-symbolic subsystems. Most quantitative claims here are tied to a checked-in benchmark artifact; where a figure is not, the text says so at the point the figure is reported. The subsystem measurements in \Cref{sec:wcoj-eval,sec:cache-eval,sec:runtime-eval} are \emph{single-system ablations of xlog against its own baselines} on development hardware; \Cref{sec:h2h} then adds controlled head-to-head comparisons against external engines (Scallop, ProbLog2, Soufflé) collected on separate cloud GPUs, \Cref{sec:ec-eval} tabulates the Event-Calculus rule-induction runs of \Cref{sec:ecinduction} beside the figures published for the same benchmark, and \Cref{sec:overhead} isolates xlog's two design costs---verification and residency---from ordinary GPU acceleration. Throughput-target envelopes for operations measured only in isolation are maintained as an internal regression contract and are not restated here, and \Cref{sec:noquant} names the capabilities that ship with correctness evidence but no timing artifact.

\subsection{Methodology}

All measurements were collected on an NVIDIA RTX PRO 3000 (Blackwell, 12\,GB VRAM, compute capability 12.0). Rust crates are built in release mode; CUDA modules are staged cubin-first with PTX fallback and explicit JIT warm-up; Python components use \texttt{pyxlog} with CUDA-enabled PyTorch. The benchmark harness uses Criterion.rs with the settings in \Cref{tab:stat-protocol}, and GPU measurements include warm-up to amortize kernel-module loading and CUDA context setup.

\begin{table}[htbp]
  \centering
  \caption{Benchmark harness settings.}
  \label{tab:stat-protocol}
  \small
  \begin{tabular}{@{}ll@{}}
    \toprule
    \textbf{Setting} & \textbf{Value} \\
    \midrule
    Sample size       & 10--100 runs (benchmark-dependent) \\
    Warm-up           & 3 iterations (PTX JIT, memory-pool init) \\
    Significance level & 0.10 \\
    Noise threshold   & 0.05 (ignore ${<}5\%$ variance) \\
    Seeding           & Deterministic LCG \\
    \bottomrule
  \end{tabular}
\end{table}

\subsection{Worst-Case Optimal Joins}\label{sec:wcoj-eval}

The WCOJ subsystem (\Cref{sec:datalog}) is exercised against the binary-join baseline by a paper-scale harness (\texttt{crates/xlog-integration/benches/}\allowbreak\texttt{wcoj\_paper\_class.rs}) over four fixtures---call-graph edges, Andersen points-to~\cite{andersen1994}, \texttt{ddisasm}-style disassembly~\cite{ddisasm}, and a recursive neural-symbolic mining analog---each with $4.19$M input rows. The scale is deliberate: WCOJ speedups depend strongly on input size and skew, and synthetic micro-benchmarks systematically overstate them, so multi-rule recursive workloads at millions of tuples are the honest setting. Both paths must produce identical row sets before any timing is taken, and each measurement is a median over 10 runs with coefficient of variation ${\le}0.05$.

\Cref{fig:wcoj} reports the speedups: a $27.96\times$ geometric mean, with every fixture between $26.6\times$ and $29.6\times$. The run that produced these figures is not among the committed artifacts. This is an ablation against xlog's \emph{own} binary-join baseline; \Cref{sec:h2h} reports the external counterpart, a fused-counting comparison against Soufflé.

\begin{figure}[htbp]
  \centering
  \includegraphics{results.pdf}
  \caption{WCOJ vs.\ binary-join speedup on four paper-scale fixtures ($4.19$M input rows each). Single-system ablation; the dashed line is the geometric mean.}
  \label{fig:wcoj}
\end{figure}

On uniform or empty inputs the adaptive classifier falls back to the binary-join chain. Peak observed allocation was about $2.3$\,GB, within the device's 12\,GB budget.

\subsection{Circuit Caching and Training}\label{sec:cache-eval}

Neural-symbolic training is measured on MNIST single-digit addition (512 training images, 5 epochs, batch size 64). Circuit caching (\Cref{sec:prob}) amortizes D4 compilation across epochs: a cache ablation over 3 seeds and 3 epochs reduced mean training time from $242.9$\,s (uncached) to $88.9$\,s (cached), a $2.74\times$ speedup (95\% CI $[2.29, 3.18]$). \Cref{tab:training-summary} summarizes the full-run timing: the first epoch pays cold-start compilation and verification, while warm epochs run only the cached forward/backward kernels. These timings were taken on the current pipeline; the run that produced them is not among the committed artifacts, which record an earlier generation of the code. The ablation is a separate experiment from this full run---its uncached arm recompiles the circuit every epoch while its cached arm compiles once, so it isolates compilation amortization; the apparent speedup therefore grows with epoch count toward the cold/warm ratio rather than being a fixed property. This 512-image protocol isolates cold/warm compilation timing and is deliberately under-trained (near-chance accuracy); held-out accuracy is measured at scale in the Scallop head-to-head (\Cref{sec:h2h}), where the same pipeline reaches ${\approx}95\%$.

\begin{table}[htbp]
  \centering
  \caption{MNIST-addition training summary. The timing rows come from the 512-image, 5-epoch training run; the circuit-cache speedup comes from a separate ablation (512 images, 3 epochs, 3 seeds) and is the mean of its three per-seed ratios, not the ratio of its means.}
  \label{tab:training-summary}
  \small
  \begin{tabularx}{\linewidth}{@{}l l L@{}}
    \toprule
    \textbf{Metric} & \textbf{Value} & \textbf{Notes} \\
    \midrule
    First epoch (cold)         & $71.7$\,s  & D4 compile + CDCL verify \\
    Steady-state epoch (warm)  & $0.25$\,s  & cached forward/backward path \\
    Total training (5 epochs)  & $72.6$\,s  & \\
    Per-query cost (steady state) & $0.97$\,ms & forward + backward; steady epoch over 256 pairs, cold epoch excluded \\
    Circuit-cache speedup      & $2.74\times$ & 95\% CI $[2.29, 3.18]$ \\
    \bottomrule
  \end{tabularx}
\end{table}

\subsection{Head-to-Head Comparisons}\label{sec:h2h}

The measurements above compare xlog against its own baselines. We additionally ran controlled head-to-head comparisons against external engines---one per reasoning mode---on ephemeral cloud GPUs (RunPod RTX 3090/4090), with matched programs, data, and metrics; the artifacts are under \texttt{artifacts/head-to-head/}. The picture is deliberately honest: xlog wins clearly where its GPU-resident, bounded-memory design is the point, and does not where a mature CPU compiler is already fast enough on small inputs. \Cref{tab:h2h} summarizes; the paragraphs below give the detail.

\begin{table*}[htbp]
  \centering
  \caption{Head-to-head summary: one external engine per reasoning mode, matched programs, data, and metrics, on ephemeral cloud GPUs. The correctness gate passes in all three; timing is the honest comparative picture. Artifacts in \texttt{artifacts/head-to-head/}.}
  \label{tab:h2h}
  \small
  \begin{tabularx}{\linewidth}{@{}l l l L@{}}
    \toprule
    \textbf{Mode / task} & \textbf{Baseline} & \textbf{Correctness} & \textbf{Timing (xlog vs.\ baseline)} \\
    \midrule
    Neural: MNIST addition           & Scallop  & acc $0.955$ vs.\ $0.950$ & steady epoch $8.71$\,s vs.\ $24.35$\,s ($\mathbf{2.80\times}$ faster); faster end-to-end \\
    Probabilistic: exact inference   & ProbLog2 & $0$ error (both) & per query ${\approx}0.25$\,s vs.\ $0.014$--$0.028$\,s ($9$--$19\times$ \emph{slower} on small programs) \\
    Deterministic: triangle counting & Soufflé  & counts match exactly & $0.24$--$0.33$\,s vs.\ $3.0$--$13.6$\,s ($\mathbf{12.6}$--$\mathbf{42.5\times}$ faster), peak memory flat $4$--$15$\,MB \\
    \bottomrule
  \end{tabularx}
\end{table*}

\textbf{Neural-symbolic: MNIST addition vs.\ Scallop.} With identical MNISTNet, data, metric, and seeds, and held-out addition accuracy on the 10k-image test set, at 20\,000 training images over 5 epochs (3 seeds) xlog and Scallop reach comparable accuracy ($0.955 \pm 0.003$ vs.\ $0.950 \pm 0.006$), while xlog runs a steady epoch in $8.71 \pm 0.09$\,s against Scallop's $24.35 \pm 0.51$\,s---a $2.80\times$ per-epoch speedup---and is faster end-to-end ($108.5$\,s vs.\ $122.7$\,s over 5 epochs) once the one-time D4 compilation amortizes. The advantage is time and residency, not accuracy.

\textbf{Probabilistic: exact inference vs.\ ProbLog2.} On five programs (a conditioned Bayesian fragment and probabilistic reach chains) xlog's exact query probabilities match the analytic answers to zero error, identical to ProbLog2 (correctness gate $|\Delta| < 10^{-4}$). On timing, however, xlog is \emph{not} competitive on these small programs: its per-query time is flat at ${\approx}0.25$\,s (fixed GPU launch plus per-call D4 compile and CDCL verify) while ProbLog2's mature CPU compiler runs each in $0.014$--$0.028$\,s ($9$--$19\times$ faster). We therefore make no exact-inference \emph{speed} claim; the probabilistic contribution is correctness-equivalent inference that is machine-verified on-device and GPU-resident, which amortizes only at scale or under repeated evaluation.

\textbf{Deterministic: WCOJ triangle counting vs.\ Soufflé.} On hub-skewed graphs, xlog's fused worst-case-optimal triangle counting produces the same triangle totals as Soufflé at every size ($5.4$M, $13.1$M, $23.1$M triangles) while running $12.6\times$, $23.9\times$, and $42.5\times$ faster on wall-clock ($0.24$--$0.33$\,s vs.\ $3.0$--$13.6$\,s), the gap widening with graph size. Crucially, fused counting never materializes the triangles: peak device memory stays flat at $4$--$15$\,MB across all sizes, whereas materializing the join intermediate needs $3$--$7$\,GB and exhausts the budget (OOM) at $23$M triangles. This is the worst-case-optimal thesis demonstrated externally---the win is bounded intermediate memory on skewed, large-intermediate workloads, not a universal Datalog speedup (on moderate skew the advantage over xlog's own binary join is only $1.1$--$1.5\times$). The three-way artifact is marked not-fully-acceptable precisely because the enumerate-then-count arm OOMs at the largest size, which is the intended demonstration of the memory bound; the fused-vs-Soufflé correctness gate passes at every size.

\subsection{Event-Calculus Rule Induction}\label{sec:ec-eval}

\Cref{sec:ecinduction} states the induction protocol and analyses its runs, and \Cref{tab:ecinduction} there lists them; this subsection adds the two things an evaluation owes. The first is the set of figures other authors have published for the same benchmark, tabulated beside our pre-registered row. The second is the runs that do not favour the system, which are reported here for the same reason the favourable one is.

\textbf{Published figures for the same benchmark.} \Cref{tab:ec-published} places our pre-registered ten-fold row beside the four published frame-level \texttt{holdsAt} F1 figures for \texttt{meeting} on the whole distributed CAVIAR corpus. Ours is $0.7329$ (precision $0.6580$, recall $0.8271$; 1516 true positives, 788 false positives, 317 false negatives), micro-averaged over pooled per-fold counts as the published evaluations are. The published figures span $0.735$ to $0.887$. Our figure is \emph{below the floor of that span}, by $0.0021$: a fully pre-registered run lands just under the lowest published figure, which is the hand-crafted rule set. We do not describe it as falling inside a parity band, because it does not fall inside one. The second-iteration run of \Cref{sec:ecinduction} reaches $0.7782$, which is inside the span, but its Event-Calculus vocabulary was chosen after the first run's test-side failure mode was known on the same folds; it is reported only with that label attached and is not the row in the table.

One asymmetry runs in the published rows' favour and we state it without drawing a conclusion from it: our gates, seeds and folds were fixed before the run that produced our row and that run was executed once, whereas the published settings are described by their authors as the best among several parameter settings they tried. We make no adjustment to any number on that basis, and we do not order the rows.

\begin{table*}[htbp]
  \centering
  \caption{Frame-level \texttt{holdsAt} F1 for the \texttt{meeting} fluent on the distributed CAVIAR corpus, micro-averaged over pooled counts. Rows are ordered alphabetically by system name and none is emphasised. They are deliberately \emph{not} ordered by F1: the five protocol deltas printed beneath the table are real, unquantified, and do not all act in the same direction, so an ordering by score would assert a comparison this evidence does not support. Dashes are figures the cited table does not report.}
  \label{tab:ec-published}
  \small
  \begin{tabularx}{\linewidth}{@{}l c c c L@{}}
    \toprule
    \textbf{System} & \textbf{P} & \textbf{R} & \textbf{F1} & \textbf{Source} \\
    \midrule
    Hand-crafted rules                      & $0.644$  & $0.855$  & $0.735$  & OLED~\cite{oled} Tab.~1(b), \texttt{EC\_crisp}; the same value appears in WOLED~\cite{woled} Tab.~2 \\
    OLED                                    & $0.678$  & $0.953$  & $0.792$  & OLED~\cite{oled} Tab.~1(b), whole dataset \\
    OLED, as re-run by other authors        & ---      & ---      & $0.782$  & WOLED~\cite{woled} Tab.~2 \\
    WOLED-ASP                               & ---      & ---      & $0.887$  & WOLED~\cite{woled} Tab.~2 \\
    xlog, EC${}+{}$inertia, pre-registered  & $0.6580$ & $0.8271$ & $0.7329$ & This work: 10-fold CV by video segment, all gates re-derived per fold (\Cref{sec:ecinduction}) \\
    \bottomrule
  \end{tabularx}

  \vspace{0.5em}
  {\footnotesize\raggedright
  \emph{Protocol deltas, which must travel with this table.}
  (i)~The published runs cross-validate over windowed interpretations with subsampled negatives; ours folds over whole video segments.
  (ii)~They learn full initiation \emph{and} termination programs; our row has a two-literal initiation clause and an empty termination theory, which by itself accounts for 788 of 2{,}304 predicted-positive frames being false positives.
  (iii)~Their rule language admits longer bodies---bottom clauses averaging some fifteen literals against our two-literal cap.
  (iv)~Their settings are reported as the best among several tried; ours were pre-registered.
  (v)~All rows compute F1 on \texttt{holdsAt} inferred through inertia, on the same corpus, with the same micro-aggregation---the one axis on which the protocols do match, and the reason the rows are placed side by side at all.\par}
\end{table*}

\textbf{A learned detector against precomputed geometry.} The result that depends on no published figure is the windowed-fold comparison with a direct per-timestep target. The search returns the same two-clause theory on every fold whether the proximity predicate is the precomputed geometric one or the learned \texttt{close\_nn}. Held-out F1 per fold is $0.9215 / 0.6804 / 0.7695$ with the precomputed predicate (mean $0.7905$) and $0.9215 / 0.7918 / 0.7695$ with the learned one (mean $0.8276$): the two modes agree to the digit on folds 1 and 3, and differ only on fold 2. A perceptual predicate that was never shown a distance label therefore stands in for hand-set geometry on this protocol without costing accuracy, which is the claim \Cref{sec:ecinduction} exists to make.

\textbf{The same detector under Event-Calculus cross-validation: a negative result.} Substituting \texttt{close\_nn} for the precomputed predicate \emph{inside} the Event-Calculus initiation search does not survive the ten-fold protocol: precision $0.1250$, recall $0.0005$, F1 $0.0011$ (1 true positive, 7 false positives, 1{,}832 false negatives) in $820.1$\,s of wall clock. Nine of the ten folds reject the candidate for insufficient new coverage; \Cref{sec:ecinduction} gives the diagnosis, which is a scale mismatch in the selection rule rather than a failure of the network. The row is reported rather than omitted: an evaluation that carried only the protocol-matched row would be publishing a selection of its own runs.

\textbf{Abstention under the clean protocol.} Re-run on the deduplicated, scene-family-grouped corpus, the \texttt{meeting} fluent returns no theory at all---the Event-Calculus route and the direct route both abstain on every fold (0/120/1812 and 0/106/1812 tp/fp/fn respectively), F1 $0.0$ on both. On \texttt{moving} the Event-Calculus route also abstains (0/0/3136), while the direct route selects a clause on nine of the ten folds---fold 6 abstains---for micro precision $0.5334$, recall $0.3817$, F1 $0.4450$ (1197/1047/1939). These zeros are empty theories returned by the permutation-null fit gate, not scored predictions: with eleven observable \texttt{meeting} initiations in the clean corpus the gate declines to commit a theory, which is the behaviour it was pre-registered to have.

\textbf{Why the two corpora differ.} A frame-level audit of every \texttt{meeting} transition event in the distributed corpus against the original ground-truth XML classifies 21 of 25 events as real, 3 as duplicates and 1 as a splice artifact, with one source video contributing 855 duplicated gold pair-frames. That defect is the reason the distributed-corpus rows carry a caveat---and also the reason the comparison in \Cref{tab:ec-published} is drawn there rather than on the clean corpus: the distributed dump is the regime the published figures share.

\subsection{Isolating Verification and Residency Overhead}\label{sec:overhead}

Two isolations separate xlog's two design costs from ordinary GPU acceleration; both run on a RunPod RTX 3090, with artifacts under \texttt{artifacts/head-to-head/}.

\textbf{Verification cost.} The cold epoch's compilation time is dominated by the on-GPU CDCL equivalence check, not the D4 compile. Splitting the cold compile via \texttt{warmup\_breakdown()} on probabilistic reachability chains, CDCL verification accounts for $96$--$100\%$ of the cold-compile time and grows with circuit size ($251$\,ms at $n{=}5$ to $508$\,ms at $n{=}40$), while D4 compilation stays under $17$\,ms. The ``expensive'' cold epoch is thus almost entirely the price of the machine-checked correctness certificate (\Cref{sec:prob})---a cost the other systems in \Cref{tab:comparison} do not pay because they do not verify.

\textbf{Residency benefit.} To isolate the transfer-free advantage from ordinary GPU acceleration we run the same neural--symbolic pipeline with and without a forced device$\to$host$\to$device round-trip at the neural--symbolic boundary---the transfer a CPU-reasoning hybrid pays every step---sweeping the number of neural$\to$symbolic handoffs per step (the batched addition path, $2$ to $512$ handoffs). Each handoff costs a stable ${\approx}40$--$56\,\mu$s round-trip; because the batched circuit compute grows sublinearly with the batch while the transfer grows linearly with the handoff count, the round-trip's share of a step \emph{rises with scale}: $0.5\%$ at $2$ handoffs, $4\%$ at $32$, and ${\approx}10\%$ at $128$---the standard batch-$64$ MNIST-addition step ($7.2$\,ms of $72$\,ms)---holding near $9\%$ at $512$. On a single query the transfer is negligible, but at a realistic training batch it is already a ${\approx}10\%$ tax that xlog avoids: a real, measured contribution, additive to and distinct from the GPU-compute and circuit-caching advantages, and growing with the per-step neural$\to$symbolic fan-out. The forced per-buffer copy is an upper bound (a hybrid could coalesce transfers), and the larger transfer-dominated regime---large-state Datalog fixpoints where whole relations would be shipped each iteration (\Cref{sec:intro})---needs a separate relation-residency ablation we leave to future work (\Cref{sec:limits}).

\subsection{Runtime Optimization}\label{sec:runtime-eval}

Two of the runtime optimizations of \Cref{sec:datalog} target repeated-session workloads. The persistent hash-index manager retains built join indices across evaluations, keyed on relation generation, schema signature and device ordinal, under a byte-budgeted LRU policy: a session whose relations change little between calls reuses an index instead of rebuilding it, and a generation change invalidates the entry rather than serving a stale one. The profile-gated shared-memory chain scorer (\texttt{crates/xlog-cuda/kernels/}\allowbreak\texttt{ilp\_exact.cu}) tiles the left relation through shared memory when scoring chain-shaped candidate bodies during exact rule induction, and is gated on a candidate-row threshold below which the baseline scorer runs unchanged.

Both are measured on repeated-session fixtures. The index manager records a $3.21\times$ speedup on a build-heavy repeated-session semi-join (cached median $0.08$\,s vs.\ uncached $0.26$\,s) with zero tracked host transfers, and the chain scorer records a $5.58\times$ speedup on a chain-hot fixture with a bounded $1.3\%$ regression on the small-input control. The runs that produced these two figures are not among the committed artifacts. They are point medians on single fixtures; per-fixture dispersion is left to future work.

What committed tests pin is narrower than these two figures. For the index manager they pin the cache bookkeeping---that the key separates entries differing in relation generation, schema or device, that the byte budget evicts by LRU, and that invalidation removes the entry---and, on a repeated-session semi-join, that the index is built once and reused thereafter with no stale acceptance, while both the cached and uncached paths record zero tracked host-to-device and device-to-host calls. That zero-transfer property is asserted both in the Rust executor tests and end-to-end through the Python session API; it is read from the provider's own transfer counters, and the two paths are compared by output row count rather than by row set. For the chain scorer no committed test executes the shared-memory kernel at all: the exact-induction tests that check the native path against a reference implementation use fixtures below the gate threshold and therefore exercise only the baseline scorer, so neither the small-input fallback nor the threshold is pinned. Common-subexpression elimination and adaptive re-optimization are validated for output equivalence rather than reported as headline speedups.

\subsection{Capability Comparison}

\Cref{tab:comparison} positions xlog qualitatively against representative systems; each addresses a subset of the design space, and xlog's contribution is unifying them in one GPU-resident runtime. This table is capability coverage, not performance: quantitative head-to-head results against Scallop, ProbLog2, and Soufflé are in \Cref{sec:h2h}, while GPUlog, VFLog, Lobster, and DeepProbLog remain to be benchmarked directly.

The rule-induction row needs its evidence status stated, because the two paths behind it are not equally evidenced and the distinction matters. The trainable-rule-body path is carried onto an external benchmark in \Cref{sec:ecinduction} and tabulated in \Cref{sec:ec-eval}, under pre-registered gates and cross-validation; it is no longer a demonstration-only surface. What that evidence does \emph{not} supply is an ordering against a dedicated rule learner, for the protocol reasons given with \Cref{tab:ec-published}. The differentiable-ILP path is not exercised there at all and remains validated by demonstrations rather than benchmarks (\Cref{sec:nesy}). Both are beta surfaces (\Cref{sec:limits}), and the table's ``yes'' claims capability, not maturity.

\begin{table*}[htbp]
  \centering
  \caption{Capability comparison across neurosymbolic and GPU logic systems.}
  \label{tab:comparison}
  \small
  \begin{tabularx}{\linewidth}{@{}l L L L L L@{}}
    \toprule
    \textbf{Dimension} & \textbf{DeepProbLog} & \textbf{Scallop} & \textbf{Lobster} & \textbf{GPUlog} & \textbf{xlog} \\
    \midrule
    Symbolic execution      & CPU (Prolog) & CPU (Datalog)    & GPU (Datalog)      & GPU (HISA) & GPU (semi-naive) \\
    Probabilistic inference & Exact (CPU)  & Provenance (CPU) & Provenance (GPU)   & None       & Exact d-DNNF + MC (GPU) \\
    Circuit verification    & None         & None             & None               & ---        & On-GPU CDCL \\
    Neural integration      & Prolog bridge & Differentiable  & Differentiable (GPU) & None     & Zero-copy, fused \\
    Zero-copy ML interop    & No           & Partial          & Partial            & No         & Yes \\
    Rule structure induction & No          & No               & No                 & No         & Yes (beta): trainable bodies, dILP \\
    Epistemic reasoning     & No           & No               & No                 & No         & Yes (finite fragment) \\
    \bottomrule
  \end{tabularx}
\end{table*}

\subsection{Capabilities Carried by Correctness Evidence Only}\label{sec:noquant}

Four capabilities added since the measured runs above ship with correctness evidence---tests that pin their behaviour---but with no committed timing artifact. We describe what each does and publish no speedup for any of them.

\textbf{Exact log-evidence and the CNF-variable-to-fact map.} An exact evaluation now returns the log partition function of the weighted formula, $\log Z$, alongside the per-query probabilities, so a program's evidence mass is read off directly instead of being reconstructed from marginals; it surfaces as \texttt{log\_z\_e} on the Python result. Paired with it, \texttt{prob\_var\_map()} reports which probabilistic fact each CNF variable stands for---an ordinary fact, one Bernoulli decision of an annotated disjunction's chain, or padding---so the per-variable gradient buffers can be attributed back to program facts rather than to opaque indices. The map is deliberately partial and says so: its length is the encoder's variable \emph{capacity}, not a variable count, and it is refused with a typed error for Monte Carlo programs and for exact programs compiled through the GPU count-lift fast path, which never builds a CNF encoding and therefore has no map to report.

\textbf{Same-head rule unions in one multiway pass.} A predicate head defined by many rules used to fold its per-rule contributions into the head relation one union at a time, re-sorting the growing accumulator once per rule and going quadratic in the rule count---a real cost at corpus scale, where one head can carry thousands of rules. Contributions are now merged with a single multiway union per head per iteration. The regression test pins this at the \texttt{-{}-stats} operation level, asserting that the union count does not scale with the same-head rule count (eight same-head rules record one union, not eight) while the derived rows stay byte-identical to the unbatched result.

\textbf{An exact memoized-DP stage beyond enumeration capacity.} The joint-constraint solver is exact by complete enumeration where a component fits the enumeration capacity. Chain-order path components that exceed it now take a memoized dynamic-programming CUDA stage over reached-domain bitsets, which is also exact: the passes are restricted forward passes, so every emitted total is a linearly accumulated \texttt{f32} and margins come only from exact passes, never from bounds. Components with wider frontiers refuse on the device with a typed status rather than silently degrading to an approximation; the pinned width gates eligibility, and the fuel meter reconciles against device-measured DP transitions instead of an estimate.

\textbf{The joint-constraint carrier and its zero-copy binding.} The carrier owns every solver buffer on the device. Python consumers receive DLPack~\cite{dlpack} views of the exact producer allocations---ownership stays with xlog, and the shared runtime identity is kept alive so strict launch recorders keep recording the exported view rather than rejecting it---and drive the cold-path session lifecycle explicitly: allocate, register schema, bind signatures, solve label feasibility. The fuel meter lives inside the session, so exhaustion saturates across calls and a retry repeats the identical typed refusal instead of slipping further work past the budget; every refusal names the precondition it enforced and its literals.

\subsection{CUDA Certification}

The CUDA kernel certification suite covers all GPU kernel operations: 207 tests across 33 categories (25 infrastructure C01--C25 + 8 GPU-specific G01--G08), at 100\% pass on the latest full certification, covering 22 kernel modules. The infrastructure categories cover toolchain validation, launch configuration, memory spaces, synchronization, warp-level operations, atomics, floating-point and determinism semantics, and host-device transfer accounting; the GPU-specific categories validate circuit forward/backward evaluation, weight injection, transfer efficiency, circuit caching, PTX robustness, and SAT/CDCL solving. Because the public repository does not rely on hosted GPUs, the suite runs through GPU-backed local release validation; the always-on CI covers formatting, linting, and no-GPU build checks.
